% ============================================================
\chapter{この資料の作り直し方}\label{app:rebuild}

\section{手順}

\begin{enumerate}
\item リポジトリの直下で、数値・表・図を作る（約 1〜2 分）。
\begin{quote}\small
\code{.venv\textbackslash Scripts\textbackslash python.exe 資料集\textbackslash LLM-BLMビューとリスクのコツ\textbackslash ビューとリスクのコツ\_計算.py}
\end{quote}
節の名前（A、B、C、D、E、E2、F、G、H1、H2、I、J、K、L、M、N、O、P）を引数に付けると、その節だけを計算し直す（B は A に、I〜K は H2 に、L は A に依存する）。
\item 資料のフォルダで、TeX をコンパイルする（tectonic は XeTeX 系の TeX で、必要なパッケージとフォントを自動で取ってくる）。
\begin{quote}\small
\code{tectonic BLMビューとリスクのコツ.tex}
\end{quote}
\end{enumerate}

計算は \code{.env} の \code{DB\_PATH}（\code{data/blm\_2015.sqlite}）の DB を使う。LLM の回答は実行番号 1、BL の係数（$\tau$、$\delta$、共分散の期間）は \code{.env} の値である。乱数を使う節（D のシミュレーション、I のシナリオ）は、節ごとに種を固定している。

\section{ファイル}

\begin{center}
\small
\begin{tabular}{p{5.0cm}p{8.6cm}}
\toprule
ファイル & 内容 \\
\midrule
\code{BLMビューとリスクのコツ.tex} & 本体（プリアンブルと、章のファイルの読み込み） \\
\code{章/} & 本文（章ごとのファイル） \\
\code{ビューとリスクのコツ\_計算.py} & 数値・表・図を作るスクリプト \\
\code{生成/} & スクリプトが作ったもの（\code{values.tex}、\code{tab\_*.tex}、\code{fig\_*.pdf}）。手で直さない \\
\code{データ/} & 外部のデータ（下の表）を CSV に整えたもの \\
\bottomrule
\end{tabular}
\end{center}

\section{外部のデータ}

\begin{center}
\small
\begin{tabular}{p{3.4cm}p{4.6cm}p{5.2cm}}
\toprule
ファイル & 中身 & 出どころ \\
\midrule
\code{ff\_factors\_\allowbreak monthly.csv} & Fama--French の 5 ファクターとモメンタム（月次、1963 年 7 月〜） & Kenneth French の Data Library（\code{F-F\_Research\_\allowbreak Data\_5\_\allowbreak Factors\_2x3}、\code{F-F\_Momentum\_\allowbreak Factor}） \\
\code{mom\_monthly.csv} & モメンタム・ファクター（月次、1927 年〜） & 同上 \\
\code{mom\_daily.csv} & モメンタム・ファクター（日次、1926 年 11 月〜） & 同上（\code{F-F\_Momentum\_\allowbreak Factor\_daily}） \\
\code{epu\_monthly.csv} & 経済政策不確実性指数（新聞ベース、月次、1900 年〜） & policyuncertainty.com（\code{US\_Policy\_\allowbreak Uncertainty\_\allowbreak Data.xlsx}） \\
\code{gpr\_monthly.csv} & 地政学リスク指数 GPR、GPRT、GPRA（月次、1985 年〜） & matteoiacoviello.com（\code{data\_gpr\_export.xls}） \\
\bottomrule
\end{tabular}
\end{center}

いずれも 2026 年 10 月 15 日に取得した。Kenneth French のデータは 2026 年 8 月の CRSP に基づく版である。データは更新されることがあるので、取り直すと数値が少し変わりうる。

\section{論文の数値例}

\begin{itemize}
\item Idzorek の 8 資産の例：論文（2004 年 7 月 20 日の版）の表 1（均衡リターン）、表 2（時価総額の比率）、表 5（共分散）、本文のビューと確信度。
\item Meucci の 6 か国の例：論文（2010 年 10 月 12 日の版）の標準偏差、相関、時価総額の比率、ビュー、$\lambda=1.2$、$\tau=0.4$。
\item 対象論文の誤差：第 2 版（2025 年 10 月 19 日）の表 5 の MSE。
\end{itemize}
